%==============================================================================
%  CAPÍTULO V
%==============================================================================
\chapter{Procedimiento concursal de reorganización simplificada de MIPEs}
\label{cap:reorganizacion-simplificada}

\fichaCapitulo{Reestructuración ágil para la pequeña empresa.}{Reemplazo de auditorías,
veedores especializados de Categoría B y el reintento del artículo 286 S.}

%==============================================================================
\bloqueTeorico
%==============================================================================

\subsection{La reducción de barreras procedimentales}

La reorganización simplificada, incorporada por la \Lreforma{} \citeyearpar{ley-21563}, no
es un procedimiento conceptualmente distinto del ordinario: es el mismo procedimiento
despojado de las exigencias que resultaban desproporcionadas para una empresa de menor
tamaño. Su técnica legislativa lo confirma: los artículos 286 A y siguientes remiten una y
otra vez a las reglas generales del Capítulo III, apartándose de ellas sólo donde la escala
de la MIPE lo justifica.

El diagnóstico que motivó su creación era conocido: bajo el texto original de 2014, el
costo de acceso al procedimiento ordinario ---auditoría, honorarios, complejidad
documental--- superaba con frecuencia el valor de la empresa que se pretendía salvar, de
modo que la micro y pequeña empresa quedaba de hecho excluida del salvataje y abocada a la
liquidación. Sobre ese diagnóstico y sus limitaciones persistentes, véanse
\textcite{aj-ley21563}, \textcite{prensa-df-brechas-mypes} y \textcite{prensa-exante-lavin}.

\subsection{Sustitución de la auditoría externa por declaración jurada}

La adaptación más significativa opera en la fase de entrada. El \art{286 A} dispone que los
antecedentes del \art{56} deben acompañarse, pero \destacar{el que en el procedimiento
ordinario requeriría certificación de un auditor independiente se informa dentro de la misma
declaración jurada} del deudor.

La justificación es de eficiencia: exigir auditoría a una microempresa cuya contabilidad es
elemental impone un costo que no se corresponde con la información que agrega. El informe de un
auditor externo registrado ante la \cmf{} representa un desembolso del orden de
\destacar{35 a 45 UF} para una empresa de mediana envergadura ---y considerablemente mayor en
estructuras complejas---, cifra significativa para el presupuesto de una microempresa en crisis y
muchas veces desproporcionada respecto de la simplicidad de su contabilidad. La contrapartida es un
traslado de riesgo. Lo que un auditor habría filtrado antes de presentar, ahora lo controlará ---o lo
objetará--- un acreedor durante la tramitación, o el tribunal en caso de liquidación posterior
por la vía del \art{169 A}.

\begin{foco}[La declaración jurada abarata, pero no exime de rigor]
Sustituir la auditoría por la declaración jurada del propio deudor reduce el costo, no la
exigencia de exactitud. La verificación documental previa (capítulo \ref{cap:entrevista})
deja de ser una recomendación y pasa a ser la única salvaguarda frente a la objeción y al
incidente de mala fe.
\end{foco}

%==============================================================================
\bloqueNormativo
%==============================================================================

\subsection{Presupuestos de admisibilidad y umbrales de la Ley \Nro 20.416}

El procedimiento se reserva a las empresas deudoras que califican como micro o pequeñas
conforme a los umbrales de ventas de la \Lpyme{} \citeyearpar{ley-20416} ---hasta
\uf{25.000} de ventas anuales--- según la sistematización del capítulo
\ref{cap:estatuto-deudores}. La calificación se acredita con la carpeta tributaria del
\sii{}.

\subsection{Nominación del veedor y antecedentes del artículo 286 A}

El \art{286 A} regula la nominación: el deudor presenta a la \superir{} copia del documento
del \art{54}, con el cargo del tribunal, y acompaña los antecedentes del \art{56}, con la
particularidad ya señalada de que el antecedente que exigiría auditoría se informa en la
propia declaración jurada.

El modelo de propuesta de acuerdo y la declaración jurada aplicables se rigen,
respectivamente, por las \ncg{24} \citeyearpar{ncg-24} y \ncg{23} \citeyearpar{ncg-23}.

\subsection{La Resolución de Reorganización simplificada}

El \art{286 B} dispone que, dentro del quinto día de la presentación, el tribunal dicta la
resolución que designa al veedor y otorga una \pfc{} \destacar{por cuarenta días} contados
desde la notificación, prorrogable conforme al \art{286 C}, \destacar{en los mismos
términos del \art{57}}. Rige, pues, la misma inmunidad ejecutiva y la misma excepción
laboral analizadas en el capítulo \ref{cap:reorganizacion-ordinaria}.

\subsection{Prórroga de la protección: el veto del 70\,\%}

El \art{286 C} regula una prórroga de hasta \destacar{treinta días}, solicitada ante el
tribunal y publicada en el \boletin{} hasta el décimo día anterior al vencimiento. Los
acreedores tienen tres días para oponerse. Vencido ese plazo, \destacar{el tribunal debe
acoger la solicitud, salvo que uno o más acreedores que representen más del 70\,\% del
pasivo se opongan}.

\begin{foco}[La prórroga simplificada se concede casi de oficio]
A diferencia del procedimiento ordinario ---donde la prórroga exige reunir apoyos
positivos---, aquí la prórroga se concede salvo que se articule un veto cualificado del
70\,\% del pasivo. La carga se invierte: no es el deudor quien debe juntar apoyos, sino los
acreedores hostiles quienes deben alcanzar el 70\,\% para impedirla. Es una diferencia
estratégica de primer orden.
\end{foco}

\subsection{Posposición de personas relacionadas y bienes garantizados}

El \art{286 E} reproduce, para el procedimiento simplificado, la posposición de los créditos
de personas relacionadas no documentados noventa días antes del inicio (véase el capítulo
\ref{cap:reorganizacion-ordinaria}).

El \art{286 Q} permite al acreedor con prenda o hipoteca solicitar, dentro de los quince
días siguientes a la publicación de la Resolución de Reorganización, que el tribunal declare
que el bien garantizado \destacar{no es esencial} para el giro. El tribunal puede pedir al
veedor un informe sobre la esencialidad y el avalúo. La calificación es decisiva: conforme
al \art{286 R}, sólo respecto de los bienes declarados esenciales se aplican los términos
del acuerdo a las obligaciones garantizadas.

\subsection{El rechazo del acuerdo y el artículo 286 S}

\begin{alerta}[Rectificación de un error difundido]
Es frecuente encontrar en la literatura de divulgación la afirmación de que el
\art{286 S} consagraría un «derecho de reintento» del deudor, esto es, la facultad de
proponer un nuevo acuerdo dentro de quinto día con el apoyo de más del 50\,\% del pasivo.
\destacar{El precepto no dice eso.}
\end{alerta}

El \art{286 S} se titula «Rechazo del Acuerdo» y regula la consecuencia del fracaso de la
propuesta. Su contenido efectivo es el siguiente: si la propuesta es rechazada por no
haberse obtenido el quórum necesario, o porque el deudor no otorgó su consentimiento, y
no estuviere constituida la junta de acreedores, el tribunal \destacar{deberá requerir a
la \superir{} la nominación del liquidador} conforme al \art{37}, dentro de los cinco
días siguientes a esa actuación, acompañando los antecedentes de los tres principales
acreedores que consten en la nómina de créditos reconocidos, excluidas las personas
relacionadas con el deudor. Recibido el certificado de nominación, el tribunal dicta la
resolución de liquidación sin más trámite.

El plazo de cinco días es, por tanto, el que tiene \emph{el tribunal} para requerir la
nominación, y no un plazo concedido al deudor para insistir.

\begin{foco}[Qué significa esto en la estrategia]
El rechazo de la propuesta simplificada conduce a la liquidación de modo prácticamente
automático y sin trámite intermedio. No existe una segunda oportunidad legal. Toda la
negociación debe estar cerrada \emph{antes} de la votación, porque después no hay
mecanismo de recuperación.
\end{foco}

\subsection{Enajenación de activos no esenciales}

La calificación de un bien como \destacar{no esencial} para el giro ---por la vía del \art{286 Q},
a solicitud del acreedor con garantía real, dentro de los quince días de la Resolución de
Reorganización--- tiene una consecuencia directa: el bien declarado no esencial puede
enajenarse para generar liquidez sin comprometer la continuidad operativa, y respecto de él el
acreedor garantizado no queda sujeto sin más a los términos del acuerdo (\art{286 R}).

\begin{foco}[Vender lo no esencial para financiar el acuerdo]
Distinguir el activo esencial del prescindible es una decisión estratégica de la MIPE en
reorganización: el bien no esencial se convierte en fuente de caja ---o en objeto de negociación
con el acreedor garantizado--- sin poner en riesgo la operación. La calificación temprana,
apoyada en el informe del veedor, ordena esa distinción antes de que la disputa la fije el
tribunal.
\end{foco}

%==============================================================================
\bloquePractico
%==============================================================================

\subsection{Modelo de propuesta de acuerdo de reorganización simplificada}

\begin{escrito}[Modelo — Propuesta de acuerdo de reorganización simplificada]
\small
\noindent\textsc{S.J.L. en lo Civil de} [\textsc{ciudad}] --- Causa Rol [\textsc{rol}]

\medskip
\noindent [\textsc{nombre del representante legal}], en representación de
[\textsc{razón social de la mipe}], en los autos de reorganización simplificada
caratulados «[\textsc{carátula}]», a US.\ respetuosamente digo:

\medskip
Que, en conformidad a lo dispuesto en la \LIR{}, modificada por la \Lreforma{}, vengo en
presentar formalmente la \textbf{propuesta de acuerdo de reorganización simplificada},
para someterla a la deliberación y votación de los acreedores, sobre la base de las
siguientes estipulaciones:

\begin{enumerate}[leftmargin=1.6em]
  \item \textbf{Remisión y condonación.} Se propone la condonación total de los intereses
  corrientes, moratorios, multas y comisiones devengadas a la fecha de la resolución de
  reorganización.

  \item \textbf{Plazo de gracia.} Se establece un plazo de gracia de doce meses para el
  pago de la primera cuota del capital reestructurado, contado desde la publicación del
  acuerdo firme en el \boletin{}.

  \item \textbf{Plan de amortización.} El capital insoluto verificado se pagará en 48
  cuotas mensuales, iguales y sucesivas, sin recargos, a partir del mes decimotercero.

  \item \textbf{Supervisión del acuerdo.} Se propone la designación de un interventor de
  la nómina de la \superir{}, Categoría B, con atribuciones de inspección mensual de
  flujos.
\end{enumerate}

\medskip
\noindent\textbf{Por tanto}, ruego a US.\ tener por acompañada la propuesta y ordenar su
publicación en el \boletin{}.
\end{escrito}

\begin{alerta}[Adaptación obligatoria del modelo]
Este modelo es un esqueleto. Los porcentajes, plazos y garantías deben ajustarse al flujo
proyectado del deudor y al mapa de mayorías. Una propuesta estándar presentada sin
negociación previa con los acreedores mayoritarios tiene alta probabilidad de rechazo.
\end{alerta}

%==============================================================================
\bloqueJurisprudencial
%==============================================================================

\subsection{Incidente de oposición al acuerdo simplificado}
\pendiente{Sistematizar los primeros fallos posteriores a la entrada en vigencia de la
Ley 21.563.}

\subsection{Efecto del rechazo por falta de quórum}

¿Conduce obligatoriamente a la resolución de liquidación simplificada dictada de oficio,
o cabe privilegiar la prórroga judicial para explorar la venta ordenada de la unidad
económica? \pendiente{Sistematizar criterios.}
